% FORMALIZACION_NUEVA de la elevacion de N39, compatible con (rho9,q9),
% T_d y Upd_t; arquitectura de residuo/cociente y orientacion preexistente.
% Fragmento de investigacion separado. No selecciona las reglas CKM N42.
\subsubsection{Transport de retenue arithmétique et holonomie entière du croisement arithmétique}
\label{subsec:transporte-entero-cruce}

La lecture modulaire du croisement admet une réalisation entière qui maintient
l’information accumulée lors du franchissement du bord. On développe ici cette
réalisation au moyen des quotients de transport. Son domaine est une coordinatisation
relevée du support APP ; l’assignation ultérieure d’une région à un secteur
angulaire conserve sa propre application d’incidence.

\paragraph{Coordinatisation nonadique et transport relevé.}
On emploie la décomposition exacte
\[
 \rho_9(n)=1+((n-1)\bmod9),\qquad
 q_9(n)=\left\lfloor\frac{n-1}{9}\right\rfloor,
 \qquad n=\rho_9(n)+9q_9(n).
\]
Une position relevée est
$z=(x,y,q_x,q_y)\in D_9^2\times\mathbb Z^2$, où
$X=x+9q_x$ et $Y=y+9q_y$. L'incrément entier $u$ de la première
coordonnée agit par
\[
 T^x_u(x,y,q_x,q_y)=
 \bigl(\rho_9(x+u),y,q_x+c_9(x;u),q_y\bigr),\qquad
 c_9(x;u)=\left\lfloor\frac{x+u-1}{9}\right\rfloor.
\]
La deuxième coordonnée se transporte de la même manière. Les avances
cardinales correspondent à $u=\pm1$ ; l'incrément nul conserve l'état.
Le cocycle satisfait
\[
 c_9(x;u+v)=c_9(x;u)+c_9(\rho_9(x+u);v).
\]
En effet, les deux membres sont le quotient de l'unique décomposition de
$x+u+v$ avec résidu dans $D_9$. Par conséquent,
$T^x_vT^x_u=T^x_{u+v}$ et $T^x_{-u}=(T^x_u)^{-1}$ ; les transports
des deux coordonnées commutent. L'état conserve le nombre de franchissements
du bord, même lorsque la position visible revient.

\paragraph{Feuillets relevés et mise à jour des registres.}
La réalisation entière est fixée en prolongeant les opérations arithmétiques
APP aux coordonnées transportées :
\begin{align*}
 \widetilde S(z)&=X+Y=x+y+9(q_x+q_y),\\
 \widetilde P(z)&=XY
 =xy+9(q_xy+q_yx)+81q_xq_y.
\end{align*}
À chaque position, le registre des deux feuillets est conservé :
\[
 \bigl(\rho_9(\widetilde S),q_9(\widetilde S),
       \rho_9(\widetilde P),q_9(\widetilde P)\bigr).
\]
Si $e:z\to z'$ est un lien et $T$ le feuillet sélectionné, son incrément
intégral s'obtient à partir du registre avant et après le transport :
\begin{equation}
 \delta_e\widetilde T=
 \rho_9(\widetilde T(z'))-\rho_9(\widetilde T(z))
 +9\bigl(q_9(\widetilde T(z'))-q_9(\widetilde T(z))\bigr).
 \label{eq:cruce-incremento-registro}
\end{equation}
Cette identité explicite la conservation : le saut du résidu et
l’accroissement du quotient reconstruisent ensemble la variation du feuillet.
La sélection de feuillet, le transport cardinal et la mise à jour du
registre sont des opérations distinctes. Dans cette coordinatisation, on compose d’abord la
sélection ordonnée de feuillets $(T_x,T_y)$ avec $T_d$ et avec la mise à jour
de leurs quotients ; la formule ne postule pas que tout le TPK soit réduit
à cette représentation locale.

\begin{proposition}[Courbure entière et transport de frontière]
\label{prop:cruce-stokes-integral}
Définissons
$\widetilde A_x=\delta_x\widetilde T_x$ et
$\widetilde A_y=\delta_y\widetilde T_y$, au moyen de
\eqref{eq:cruce-incremento-registro}, et soit
\[
 \widetilde F_{xy}(X,Y)=
 \widetilde A_x(X,Y)+\widetilde A_y(X+1,Y)
 -\widetilde A_x(X,Y+1)-\widetilde A_y(X,Y).
\]
Les choix homogènes ont une courbure nulle. Les choix mixtes
$(\widetilde S,\widetilde P)$ et
$(\widetilde P,\widetilde S)$ ont respectivement pour courbures les entiers
$+1$ et $-1$. Pour le bord positif d'un rectangle
$R_{a,b}$ de côtés entiers $a,b>0$,
\[
 \widetilde W_{SP}(\partial R_{a,b})=ab,\qquad
 \widetilde W_{PS}(\partial R_{a,b})=-ab.
\]
Le résultat est indépendant du point de base et du nombre de franchissements
des raccords nonadiques.
\end{proposition}

\begin{proof}
La reconstruction exacte des deux feuillets donne
\[
 \delta_x\widetilde S=\delta_y\widetilde S=1,
 \qquad\delta_x\widetilde P=Y,
 \qquad\delta_y\widetilde P=X.
\]
Ainsi, la connexion $SP$ est $(1,X)$, de courbure
$1+(X+1)-1-X=1$ ; la connexion $PS$ est $(Y,1)$, de courbure
$Y+1-(Y+1)-1=-1$. Pour un feuillet unique, les différences
mixtes s'annulent. La somme des courbures des $ab$ plaquettes annule chaque
lien intérieur avec son inverse et conserve les liens de frontière.
On obtient ainsi $\pm ab$ dans $\mathbb Z$. Les quotients transportés
garantissent que la même identité vaut lors du franchissement de tout raccord.
\end{proof}

L'accumulateur de frontière peut à son tour être conservé sous forme nonadique :
$W=r_W+9q_W$. Pour une contribution entière $w_e$, il est mis à jour par
\[
 r_W'=\rho_9(r_W+w_e),\qquad
 q_W'=q_W+\left\lfloor\frac{r_W+w_e-1}{9}\right\rfloor.
\]
La composition par liens fournit la même somme que son évaluation
entière directe. Dans cette coordinatisation, le zéro entier a pour coordonnées
$(r_W,q_W)=(9,-1)$ ; c’est une représentation arithmétique de l’accumulateur,
sans identification supplémentaire avec les différents zéros orientés du TRIT.

\paragraph{Trois opérations sur l'orientation.}
Soit $\gamma$ une route cardinale fermée. Son inversion temporelle
$\gamma^{-1}$ inverse l'ordre des liens et l'orientation de chacun ;
par additivité,
\[
 \widetilde W_{SP}(\gamma^{-1})=-\widetilde W_{SP}(\gamma).
\]
L'échange des feuillets satisfait également
\[
 \widetilde W_{PS}(\gamma)=-\widetilde W_{SP}(\gamma).
\]
Pour le vérifier sur toute route, observons que la somme des deux
connexions est l'incrément exact de
$\widetilde S+\widetilde P=X+Y+XY$ ; son intégrale sur une boucle est nulle.

Le demi-tour simultané des deux directions agit, d’autre part,
par $J_c(X,Y)=(2c_x-X,2c_y-Y)$, avec $2c\in\mathbb Z^2$,
de sorte qu’il préserve la coordinatisation du réseau. Pour un lacet, on a
\[
 \widetilde W_{SP}(J_c\gamma)=\widetilde W_{SP}(\gamma).
\]
En effet, la contribution horizontale fermée est nulle et la verticale
se transforme comme
$\sum (2c_x-X)(-\delta Y)=\sum X\delta Y$, car
$\sum\delta Y=0$. Un demi-tour spatial préserve donc
l’orientation aréale de cette coordinatisation. Sa combinaison avec l’inversion temporelle
ou avec l’échange de feuillets se calcule en composant les opérations
correspondantes.

La différence est pertinente pour le registre électronique : l'inversion
simultanée des deux déplacements tous les $27$ pas d'avancement produit, dans
des fenêtres de $9$, la suite
\[
 \tau_j=(-1)^{\lfloor j/3\rfloor},\qquad 0\leq j<12,
 \qquad (+,+,+,-,-,-,+,+,+,-,-,-).
\]
Cette suite enregistre les orientations d'avancement. L'holonomie requiert
en outre les liens parcourus, les feuillets sélectionnés et les quotients
accumulés. Ainsi, le transport de l'orientation électronique est conservé
sans identifier un demi-tour des deux axes à l'inversion temporelle
d'une boucle.

\paragraph{Application au croisement marqué et portée sectorielle.}
Pour le croisement rectangulaire de côtés $4$ et $7$, cette réalisation évalue
exactement
\[
 \widetilde W_{SP}=28=1+9\cdot3,\qquad
 \widetilde W_{SP}(\gamma^{-1})=-28=8+9(-4).
\]
L’entier conserve l’aire orientée ; le résidu isolé représente seulement
sa classe nonadique. Par exemple, à partir de $(7,1)$, les représentants
périodiques des liens ont pour somme $-35$, tandis que la correction de
retenue ajoute $63$ et retrouve $28$.

Dans l'incidence marquée du lecteur angulaire, le $4$ est
$|H_e\cap P_\pi|$ et le $7$ est le pivot marqué $p_\varphi$.
La lecture rectangulaire fournit une réalisation entière du produit
$|H_e\cap P_\pi|\,p_\varphi$ une fois ce croisement déclaré.
L'application qui identifie cette région au secteur $12$ doit conserver
la marque et l'orientation du système sectoriel de
la Définition~\ref{def:ii09-sistema-sectorial}. Les affectations
sectorielles $12,23,13$ et leurs combinaisons de $A,h,\Delta^\circ$
sont une spécification supplémentaire à la loi d'holonomie rectangulaire.
Le présent résultat matérialise le relèvement d'aire, la retenue arithmétique et
leurs involutions, avec ces données de départ explicites.
